\documentclass[11pt]{article}

\usepackage[a4paper,margin=1in]{geometry}
\usepackage[T1]{fontenc}
\usepackage[utf8]{inputenc}
\usepackage{lmodern}
\usepackage{microtype}
\usepackage{xcolor}
\usepackage{amsmath,amssymb,amsthm}
\usepackage{enumitem}
\usepackage{fancyhdr}
\usepackage[hidelinks]{hyperref}

\definecolor{ink}{HTML}{243447}
\definecolor{accent}{HTML}{0B7285}
\setlist{itemsep=2pt,topsep=4pt}
\setlength{\parindent}{0pt}
\setlength{\parskip}{5pt}

\newcommand{\studentname}{Your Name}
\newcommand{\coursename}{Data Structures}
\newcommand{\termname}{Spring 2026}
\newcommand{\assignmentnumber}{04}

\pagestyle{fancy}
\fancyhf{}
\lhead{\textsc{\coursename}}
\chead{Assignment \assignmentnumber}
\rhead{\studentname}
\cfoot{\thepage}
\renewcommand{\headrulewidth}{0.4pt}
\renewcommand{\headrule}{\hbox to\headwidth{\color{accent}\leaders\hrule height \headrulewidth\hfill}}

\newcounter{problem}
\newenvironment{problem}[1][]
  {\refstepcounter{problem}\par\vspace{8pt}\textcolor{accent}{\large\bfseries Problem \theproblem}\if\relax\detokenize{#1}\relax\else\ --- #1\fi\par\smallskip\ignorespaces}
  {\par\medskip}
\newenvironment{solution}
  {\begingroup\color{ink}\textbf{Solution.}\quad\ignorespaces}
  {\endgroup\par}
\newtheorem*{claim}{Claim}

\title{\vspace{-2em}\textcolor{ink}{Assignment \assignmentnumber}\\
  \large A practical study of graph algorithms}
\author{\studentname\\\coursename\ --- \termname}
\date{Due 18 March 2026}

\begin{document}
\maketitle
\thispagestyle{fancy}

\textbf{Submission checklist.} Show enough work to make your reasoning
verifiable. State the running time of every algorithm and cite any external
source that informed your solution.

\begin{problem}[Breadth-first search]
Let $G=(V,E)$ be an undirected graph with $|V|=n$ vertices and $|E|=m$
edges. Describe an algorithm that returns the shortest unweighted distance
from a source $s$ to every reachable vertex. Give a loop invariant for the
queue-processing loop and prove its running time.

\begin{solution}
Place $s$ in a FIFO queue with distance zero. Whenever a vertex is removed,
inspect its adjacency list and enqueue each undiscovered neighbour with
distance one greater than the current vertex. The invariant is that a vertex
leaves the queue only after its stored distance is its shortest path length:
the first path reaching it has the smallest possible number of edges because
the queue processes vertices in nondecreasing distance order. Each vertex and
edge is inspected at most once, so the time is $O(n+m)$ and the storage is
$O(n)$.
\end{solution}
\end{problem}

\begin{problem}[A weighted counterexample]
Give a small weighted graph showing why breadth-first search cannot replace
Dijkstra's algorithm when edge weights are positive but unequal. Explain the
incorrect choice it makes.

\begin{solution}
Use vertices $s,a,b,t$ and edges
\[
  s\mathbin{-}a\ (1),\qquad a\mathbin{-}t\ (10),\qquad
  s\mathbin{-}b\ (4),\qquad b\mathbin{-}t\ (1).
\]
The two-edge path through $a$ is discovered first, but costs $11$; the path
through $b$ costs $5$. Breadth-first search minimises the number of edges,
not their total weight.
\end{solution}
\end{problem}

\begin{problem}[Design question]
A graph is stored as an adjacency matrix. What is the asymptotic cost of
testing whether an edge $(u,v)$ exists, and what is the cost of enumerating
all neighbours of $u$? Contrast both answers with an adjacency-list graph.
\end{problem}

\begin{claim}
For sparse graphs, adjacency lists use less space and make neighbour
enumeration proportional to the degree of the selected vertex \cite{cormen2009}.
\end{claim}

\bibliographystyle{plain}
\bibliography{references}
\end{document}
